\paragraph{Data.}
Online conformal methods provide many interval forecasters with the same
coverage target, which differ mainly in how they track the tails; they are
a natural test bed for comparisons at small $\alpha$. We use seven public
series: the oil temperature of an electricity
transformer, hourly (ETTh1) and every 15 minutes (ETTm1)
\citep{zhou2021informer}; the occupancy rate of one road sensor, hourly, and
one daily exchange rate \citep{lai2018modeling}; German--Luxembourg
day-ahead electricity prices, hourly, from Open Power System Data
\citep{wiese2019open}; and the first twenty series of the M4 hourly
\citep{makridakis2020m4} and NN5 daily \citep{godahewa2021monash}
collections, each concatenated into one series. After 50 warm-up rounds,
which initialise the methods and are not evaluated, the evaluation periods
have between 5{,}538 and 66{,}146 rounds; they are shorter at the 99\%
level, for the reason given below.

\paragraph{Forecasters.}
All methods share one point forecast and one scale. The point forecast
$\mu_t$ for $y_t$ is the mean of five simple forecasters (naive, seasonal
naive, rolling mean, simple exponential smoothing and a linear trend), each
using data up to round $t-h$, where $h$ is the seasonal period (24 for
hourly series, 96 for ETTm1, 7 for daily series). The scale $\sigma_t$ is
the standard deviation of the last 50 observations, floored at 5\% of the
series' standard deviation. Each method turns the standardized scores
$|y_s-\mu_s|/\sigma_s$, $s<t$, into an interval
$[\mu_t-\sigma_tq^-_t,\mu_t+\sigma_tq^+_t]$, so the intervals for round $t$
depend on information up to round $t-1$ only and satisfy the conditions of
Section~\ref{sec:setup}. The six methods are: adaptive conformal inference
aggregated over step sizes, AgACI \citep{zaffran2022adaptive}; dynamically
tuned adaptive conformal inference, DtACI \citep{gibbs2024conformal};
strongly adaptive online conformal prediction, SAOCP
\citep{bhatnagar2023improved}; quantile tracking (QT) and conformal PI
control (PI) \citep{angelopoulos2023conformal}; and AgACI-W, introduced
here, which aggregates adaptive conformal experts over several reference windows as
well as step sizes, with the ML-Poly rule of \citet{gaillard2014second}. For the methods
with a tuning parameter (the reference window of AgACI and DtACI, the
learning rate of QT and PI) we keep the configuration with the lowest mean
score on the evaluation period. The comparisons are therefore between the
best members of each family, selected in sample, and we use them to compare
the procedures rather than to rank the methods. At the 99\% level the
candidate reference windows are longer, and because they need more
history, the evaluation period of six series is 1{,}100 to 1{,}500 rounds
shorter at that level.

\paragraph{Procedures.}
For each data set and level we compare all 15 pairs of methods. The
Diebold--Mariano test is two-sided at 5\% with a Newey--West variance and
lag $\max\{h-1,\lfloor4(T/100)^{2/9}\rfloor\}$. The confidence sequences use
$\delta=0.05$. The EB sequence is the tuned one of Section~\ref{sec:cs},
tuned on the first 10\% of the evaluation period. The Hoeffding sequence
uses the uniform grid with $\lambda_0$ equal to the inverse of twice the
median range over the 50 warm-up rounds that precede the evaluation
period. We report whether each sequence excludes zero at
the last round and, for every procedure, which method it favours. As a check that uses
no model, we also split each evaluation period in halves, apply each
procedure to the first half (the EB sequence tuned on the first 10\% of
that half), and record whether the sign of its significant
results agrees with the mean score difference on the second half.

\paragraph{Results.}
Table~\ref{tab:empirical} summarizes the 105 comparisons at each level.
Misses, rounds on which at least one of the two intervals misses, make up
about $1.3\alpha$ of all rounds,
so at the 99\% level the median pair has 179 of them out of more than ten
thousand. The share of
pairs that each procedure separates falls with $\alpha$, fastest for the
range-calibrated Hoeffding sequence, which separates 44\% of pairs at 80\%
and none at 99\%, as Proposition~\ref{prop:hoeffding} predicts. The EB
sequence separates 70\% at 80\% and 18\% at 99\%, and the Diebold--Mariano
test declares 82\% and 54\% significant; the two never disagree in
direction. The untuned EB sequence, with $\lambda_0$ from the 50 warm-up
rounds as for the Hoeffding sequence, separates 59\% and 13\%.

The conditions of Section~\ref{sec:rates} roughly hold. Realized miss
rates have a median of $0.98\alpha$ to $1.00\alpha$ at every level (range
$0.6\alpha$ to $4.1\alpha$), and over windows of $20/\alpha$ rounds their
90th percentile has a median of about $1.1\alpha$ and is at most
$5.1\alpha$, so the conditional constant $C$ is small. Endpoint distances,
scaled by $\sigma_t$, have a 90th percentile about three times the median.
The effect size $\Delta/H$ has a median of 0.35, 0.37, 0.36 and 0.50 at
the four levels: it does not shrink with $\alpha$, so the $1/\alpha$ growth
of the required number of rounds applies to these data.

The split-sample check shows how reliable the significant results are. Of
the
Diebold--Mariano results that are significant on the first half, the share
whose sign agrees with the second-half mean falls from 100\% at 80\% to
97\%, 95\% and 90\% at 99\%. The sign of all 149 significant first-half
results of the EB sequence agrees, although at 99\% it has 15 such results
against 49 for the
Diebold--Mariano test. Configurations were selected on the whole period, so
the check is not fully out of sample.

\begin{table}[t]
\centering\small
\caption{Comparisons of six online conformal methods on seven data sets,
15 pairs each. Misses: rounds on which at least one interval misses,
median per pair. Significant: the
procedure separates the pair over the evaluation period (DM: two-sided
5\%; EB: tuned 95\% confidence sequence; Hoeffding: 95\% Hoeffding-type
sequence). First half:
number of significant results on the first half of the evaluation period,
and the share of those whose sign agrees with the mean on the second
half.}
\label{tab:empirical}
\begin{tabular}{rrrrrrrrr}
\toprule
 & & \multicolumn{3}{c}{significant (\%)} & \multicolumn{4}{c}{first half}\\
\cmidrule(lr){3-5}\cmidrule(lr){6-9}
 & & & & & \multicolumn{2}{c}{DM} & \multicolumn{2}{c}{EB}\\
\cmidrule(lr){6-7}\cmidrule(lr){8-9}
level & misses & DM & EB & Hoeffding & signif. & sign agrees (\%) & signif. & sign agrees (\%)\\
\midrule
80\% & 3905 & 81.9 & 69.5 & 43.8 & 72 & 100.0 & 58 & 100.0\\
90\% & 1982 & 66.7 & 47.6 & 26.7 & 66 & 97.0 & 40 & 100.0\\
95\% & 978 & 64.8 & 39.0 & 17.1 & 59 & 94.9 & 36 & 100.0\\
99\% & 179 & 54.3 & 18.1 & 0.0 & 49 & 89.8 & 15 & 100.0\\
\bottomrule
\end{tabular}

\end{table}
